\part*{PPD2, Merkblatt zum Harvard-Konzept\\\normalsize
Severin Henzi, Simon Kronenberg, Malte Scheurer, Cyril Wendl}

\section*{Das Harvard-Konzept}
Im Buch ``Getting to Yes'' (Deutsch: ``Das Harvard-Konzept'' \cite{harvardKonzeptDE}) beginnen die Autoren \textcite{Fisher1981GettingTY} mit der Frage, welches die beste Art sei, wie Menschen mit ihren unterschiedlichen Arten umgehen können und der Feststellung, dass der Alltag von Familien, Paaren, Angestellten, Vorgesetzten, Konsumenten, Verkäufern und anderen Menschen, die in gemeinsamer Beziehung stehen, von der Herausforderung geprägt ist, in vielen Bereichen Einigungen erzielen zu können (``getting to yes''), ohne sich in Streitigkeiten zu verwickeln. Dabei wird häufig einer von zwei Ansätzen gewählt, um Einigungen zu erzielen: der ``harte'' Ansatz, bei dem jede Verhandlung zum Willens- und Überlebenskampf wird, oder der ``weiche'' Ansatz, bei dem eine Person nachgibt und Zugeständnisse macht, um einen Kompromiss zu erzielen. Die Autoren plädieren jedoch für einen anderen Ansatz, das sogenannte sachbezogene Verhandeln (``\textit{principled negotiation''}), welches sie im Rahmen des Harvard Negotiation Projects erarbeitet haben \cite{hnp}. Der grundlegende Befund dieses Ansatzes ist, dass Verhandlungen nicht als Präferenzen-basiertes ``Feilschen'' auszuführen seien, sondern vielmehr sachorientiert und basierend auf möglichst fairen Standards, welche zum grösstmöglichen gemeinsamen Nutzen führen. Aufgrund der Trennung vom Sachverhalt zu den Menschen lässt sich dieser Ansatz als ``hart'' in der Sache aber ``weich'' im Umgang mit Menschen bezeichnen. 

Die Autoren stellen drei Grundanforderungen an Verhandlungen: 
\begin{enumerate*}
\item Verhandlungen sollten zu einer Übereinkunft in einem gegebenen Thema führen;
\item Verhandlungen sollten effizient sein;
\item Verhandlungen sollten das Verhältnis zwischen beiden Parteien nicht stören, sondern wenn möglich stärken.
\end{enumerate*} Darauf aufbauend entwickeln die Autoren vier Kriterien, um sachorientiertes Verhandeln zu definieren (s. \cref{table:gty})

\begin{table}[H]
\centering
\begin{tabular}{p{4cm}|p{4cm}||>{\columncolor[gray]{.9}}p{6cm}}
\multicolumn{2}{c||}{\textbf{Problemorientierte Ansätze}, ``Feilschen''} & \textbf{Lösungs-Orientierter Ansatz} \\\midrule
\textbf{Weich} & \textbf{Hart} & \textbf{Sachorientiert} \\\midrule\midrule
Teilnehmende sind Freunde & Teilnehmende sind Gegner & \textbf{1. Trennung von Person und Problem} \\\midrule
Zugeständnisse, um Beziehung zu erhalten & Forderungen durchsetzen & \textbf{2. Fokus auf Interessen, nicht Positionen} \\\midrule
Konflikte werden vermieden & Konflikte werden durchgesetzt & \textbf{3. Entwicklung von Optionen zum beidseitigen Vorteil} \\\midrule
Auf Einigung Bestehen & Auf eigene Position bestehen & \textbf{4. Bestehen auf objektive Kriterien für Lösungen} \\
\end{tabular}
\caption{Prinzipien sachorientierten Handelns nach \cite[13]{Fisher1981GettingTY}}
\label{table:gty}
\end{table}

Punkt 1 bringt zum Ausdruck, dass das Problem auf bestmögliche Weise gelöst werden müsse, unabhängig persönlicher Interessen und Vertrauensverhältnisse. Punkt 2 besagt, dass Interessen aufgedeckt und diskutiert werden müssen, statt Positionen zu debattieren oder der persönlichen Beziehung zuliebe allzu schnell eine Einigung erreichen zu wollen. Punkt 3 beinhaltet, dass verschiedene Handlungs- oder Entscheidungsoptionen zum beidseitigem Vorteil erarbeitet werden sollten, bevor direkt zur Entscheidung übergegangen wird. Punkt 4 fordert, dass ein Resultat basierend auf sachbezogenen Standards und möglichst unabhängig individueller Durchsetzungskraft erreicht werden. Dem Druck anderer Teilnehmer, die auf eine Lösung drängen (weicher Ansatz) oder die eigene Position verteidigen (harter Ansatz) soll standgehalten werden.

Für den Fall, dass keine Einigung erzielt werden kann, hilft es, sich der bestmöglichen Handlungs-Alternative bewusst zu sein. Indem man sich ein sogenanntes BATNA (\textit{Best Alternative To a Negotiated Agreement}) zurechtlegt, kann man der Neigung widerstehen, eine Verhandlung um jeden Preis abzuschliessen \cite{elementeVerhandeln}. Fällt das Ergebnis einer Verhandlung schlechter aus als das BATNA, lohnt es sich, das BATNA statt dem Resultat der Verhandlung zu wählen.

 \section*{Anwendungen im Kontext der Sekundarstufe II}
Im Kontext der Sekundarstufe II ergeben sich vielfältige Anwendungen, bei denen Verhandlungssituationen auftreten. Diese können unterteilt werden in \begin{enumerate*}
\item Verhandlungen mit Vorgesetzten (dem Rektorat);\item Verhandlungen mit Gleichgestellten (Lehrpersonen, Fachschaft);\item Verhandlungen mit SuS
\end{enumerate*}. Der erste Fall (Rektorat) wird im Video behandelt. Im Folgenden wird für die zwei weiteren Bereiche analysiert, wie gewinnbringende Verhandlungsstrategien gemäss der Harvard-Theorie in konkreten Situationen aussehen könnten.

\begin{myExBox}[title={Anwendung 1: Verhandlung mit Gleichgestellten (Fachschaft, Lehrpersonen)}]
Lehrpersonen verhandeln regelmässig mit Kolleg*innen der Fachschaft. Dabei kann es beispielsweise um die Anschaffung neuer Unterrichtsmaterialien, die Verteilung von Wochenlektionen, die Erarbeitung von Richtlinien oder um die Planung von Projektwochen gehen.

Im Unterschied zu Verhandlungen mit der Schulleitung oder mit SuS gibt es in diesem Fall kein vorbestimmtes Hierarchie-Gefälle. Eine gewisse Hierarchie ergibt sich jedoch häufig durch Erfahrung, Alter, Verantwortung (z.B. Fachschaftsvorstand) oder soziale Dynamiken. Auch hier ist es wichtig, sich bei Verhandlungen auf Interessen und nicht auf Positionen zu fokussieren: Es sollten nicht die erfahrensten Lehrpersonen der Fachschaft alles bestimmen können, sondern die Interessen aller Lehrpersonen sollen gleich gewichtet werden. Wenn es um die Organisation spannender Projektwochen geht, sollten Interessen neuer Lehrpersonen ebenfalls in die Entscheidungen einfliessen. Eine neue Lehrperson könnte dabei die Interessen der erfahreneren Kolleg*innen wertschätzen und gleichzeitig Optionen zum beidseitigen Vorteil entwickeln, um auch die eigenen Interessen geltend zu machen.  

\end{myExBox}

\begin{myExBox}[title=Anwendung 2: Verhandlung mit Schülern]
In einer Klasse diskutieren die SuS und die Lehrperson über die Verlängerung der Abgabefrist für eine Projektarbeit. Einige SuS bitten um mehr Zeit, während die Lehrperson auf eine pünktliche Abgabe besteht. Nach dem lösungsorientierten Harvard-Konzept wird in den vier oben vorgestellten Schritten vorgegangen.
Zunächst wird die Beziehung von Lehrperson und SuS vom Problem getrennt. Die Lehrperson verzichtet auf Vorwürfe wie ``Ihr seid schlecht organisiert'' und zeigt Verständnis für die Situation der SuS. Gleichzeitig respektieren die SuS die Haltung der Lehrperson, für die Verantwortung und Pünktlichkeit wichtig ist. So entsteht eine konstruktive Atmosphäre.

Im zweiten Schritt rücken die Interessen in den Fokus. Die Lehrperson möchte faire Bedingungen schaffen und den Lehrplan einhalten. Die SuS hingegen brauchen mehr Zeit aufgrund von Prüfungen in anderen Fächern oder ausserschulischer Verpflichtungen. Anstatt auf ihren jeweiligen Positionen zu beharren, arbeiten beide Seiten daran, die Beweggründe der anderen Partei zu verstehen und eine Lösung zu finden, die auch der anderen Partei gerecht wird.

Im dritten Schritt werden kreative Lösungen entwickelt. Möglichkeiten wie eine kurze Fristverlängerung für alle oder die Abgabe eines Entwurfs zum ursprünglichen Termin mit späterer Überarbeitung werden gemeinsam diskutiert. Diese Ansätze öffnen neue Wege und führen schlussendlich zu innovativen und für beide Seiten potentiell besseren Lösungen.

Während der gesamten Verhandlung sollten objektive Kriterien herangezogen werden. Prüfungstermine der SuS sowie der Lehrplan dienen als faktische Grundlage für die Verhandlung. SuS, die mehrere Prüfungen in derselben Woche haben, können die Projektarbeit mit einer kurzen Begründung später abgeben, während für alle anderen der Abgabetermin gleich bleibt. Durch diesen Ansatz finden beide Seiten eine faire und respektvolle Lösung, die auf gegenseitigem Verständnis und transparenten Kriterien beruht. Die Verhandlung endet mit einer Einigung, die beide Parteien zufriedenstellt.

\end{myExBox}




%\newpage
%\section*{Anhang}
%\begin{myExBox}[title=Fallbeschreibung Video]
%Das im Video behandelte Fallbeispiel bezieht sich auf eine echte Situation, die sich im beruflichen Umfeld der Studierenden ereignet hat.
%
%Im Kanton Zürich wird zwischen \textit{Lehrpersonen} (mit Lehrdiplom) und \textit{Lehrbeauftragten} (ohne Lehrdiplom) unterschieden. Letztere werden meistens angestellt, um Restpensen zu füllen, oder bei dringendem Bedarf nach Lehrpersonal, etwa bei krankheits- oder schwangerschaftsbedingten Ausfällen, oder wenn neue Fächer eingeführt werden, wie kürzlich das Fach Informatik, und pädagogisch ausgebildete Lehrpersonen in noch nicht ausreichender Anzahl verfügbar sind. Im Rahmen der Anstellung als Lehrbeauftragte(r) wird am Ende des ersten Jahres eine Leistungsbeurteilung erstellt aufgrund einzelner Unterrichtsbesuche und aufgrund standardisierter Umfragen, welche in den unterrichteten Klassen durchgeführt werden. Die Endbeurteilung dient dazu, um über die weitere Anstellung zu entscheiden, und fällt in vier Kategorien: \textit{Sehr gut, gut, genügend} oder \textit{ungenügend}. Im Vorliegenden Fall entschied der Rektor sich ursprünglich dazu, das Prädikat ``\textit{gut}'' zu verteilen, mit der Begründung, neue Lehrpersonen könnten \textit{per se} nicht \textit{sehr gut} sein, da ihnen die Erfahrung dazu mangelte. Der Lehrbeauftragte fühlte sich unfair behandelt, denn obgleich das Urteil keine direkten Konsequenzen hatte, hätte er sich aufgrund seines starken Engagements für das Fach Informatik mehr \textbf{Anerkennung} gewünscht. Deshalb suchte er das Gespräch mit dem Rektor, welches sich in ungefähr wie folgt abspielte.
%
%\begin{myExBox}[title=Im Studio]
%\begin{dialogue}
%\speak{Moderator} Unserer Redaktion wurde ein sehr interessantes Video zugespielt, dessen Quelle wir
%natürlich nicht verraten können. Es zeigt ein Verhandlungsgespräch zwischen einem
%Lehrbeauftragten und seinem Rektor – eine Gesprächssituation, die für Lehrpersonen von
%grossem Interesse ist.
%Zusammen mit unserem Studiogast, dem Kommunikationsexperten Herr XY, analysieren wir,
%wie Verhandlungen nach dem Harvard-Konzept erfolgreich geführt werden können.
%\speak{Moderator} Herr XY, herzlich willkommen. Es ist mir eine Ehre, Sie heute hier zu haben.
%\speak{Experte} Vielen Dank für die Einladung.
%\speak{Moderator} Als Kommunikationsexperte haben Sie viele Gespräche analysiert. Ist es für Sie neu, quasi
%hinter die Türen eines Rektoren-Zimmers zu blicken?
%\speak{Experte} Absolut. Ich bin gespannt auf die Bilder.
%\speak{Moderator} Ich auch. Dann schauen wir gleich rein...
%\end{dialogue}
%\end{myExBox}
%\textit{Donnerstagmorgen, 7:30, Büro des Rektors. Der Lehrbeauftragte (LB) trifft ein.}
%\begin{dialogue}
%\speak{Rektor} Guten Morgen! Nimm Platz. Was bringt dich zu mir?
%\speak{LB} Vielen Dank. Ich wollte über die Leistungsbeurteilung sprechen. Ich bin froh, dass mein
%Anstellungsverhältnis verlängert wurde, aber ich war überrascht, dass meine Beurteilung nur
%``gut'' und nicht ``sehr gut'' war. Können Sie mir sagen, wie ich mich verbessern kann?
%\direct{Die Kamera wechselt zum Studio}
%\speak{Experte} Das ist ein Schlüsselmoment. Der LB zeigt Interesse an den Gründen für die Bewertung, um
%darauf aufbauend eine Lösung zu entwickeln – ein zentraler Aspekt des Harvard-Konzepts. Erst
%darauf aufbauend können wir unseren weiteren Weg zu einer erfolgreichen Verhandlung in
%Angriff nehmen.
%\direct{Die Kamera wechselt zum Büro des Rektors}
%
%\speak{Rektor} Bei uns erhalten junge Lehrpersonen höchstens „gut“, da es ihnen noch an Erfahrung
%fehlt um besmögliche Unterrichtsqualität zu gewährleisten. Du machst Deine Arbeit wirklich
%gut, aber zum Beispiel bei der Rhythmisierung Deiner Unterrichtseinheiten sehe ich noch
%Verbesserungspotenzial.
%[LB]: Das verstehe ich. Mein Unterricht ist natürlich nur schwer mit erfahreren Lehrpersonen zu
%vergleichen. Allerdings habe ich mir für mein überdurchschnittliches Engagement inner- und
%ausserhalb des Unterrichts mehr Anerkennung gewünscht. Dürfte ich dir kurz die Gründe
%erklären die aus meiner Sicht für ein “sehr gut” sprechen?
%\end{dialogue}
%\end{myExBox}
\begin{myExBox}[title=Lernapplikation \& Video]
\begin{minipage}[c]{.25\textwidth}
Lernapplikation: \\
\small \texttt{\href{https://wendl.ch/quiz-harvard}{wendl.ch/quiz-harvard}}
\end{minipage}
\begin{minipage}[c]{.22\textwidth}
\qrcode[height=1.3cm]{https://wendl.ch/quiz-harvard}
\end{minipage}
\vline\hfill
\begin{minipage}[c]{.25\textwidth}
Video: \\
\small \texttt{\href{https://wendl.ch/ppl2}{wendl.ch/ppl2}}
\end{minipage}
\begin{minipage}[c]{.22\textwidth}
\qrcode[height=1.3cm]{https://wendl.ch/ppl2}
\end{minipage}
\end{myExBox}

\newpage
\printbibliography